\documentclass[aspectratio=169]{beamer}

\usepackage{fontspec}
\usepackage[math-style=ISO]{unicode-math}

\setmathfont[
  Extension=.otf,
]{STIXTwoMath-Regular}

\usepackage[style=authoryear]{biblatex}
\usepackage{amsmath}
\usepackage{amssymb}
\usepackage{proof}
\usepackage{stmaryrd}
\usepackage[absolute,overlay]{textpos}

\usetheme{default}
\usecolortheme{seahorse}
\addtobeamertemplate{navigation symbols}{}{%
    \usebeamerfont{footline}%
    \usebeamercolor[fg]{footline}%
    \hspace{1em}%
    \raisebox{1.2pt}[0pt][0pt]{\insertframenumber/\insertmainframenumber}
}

\setbeamertemplate{itemize item}[circle]
\setbeamertemplate{itemize subitem}[circle]
\setbeamertemplate{itemize subsubitem}[circle]

\title{Introduction to Proving Stuff\textsuperscript{\texttrademark} with Logical Relations}
\author{Jesse Sigal}
\date{November 7, 2024}

\addbibresource{refs.bib}

\begin{document}

\begin{frame}
\titlepage
\end{frame}

\begin{frame}
\frametitle{Attribution}
\begin{itemize}
\item Slides with $*$ at the end of their title were written with the help of GPT 4o (for lazy \LaTeX'ing).
\item Most things for the calculus are in line with \cite{Crole_1994}.
\end{itemize}
\end{frame}

\begin{frame}
\frametitle{Overview}
\begin{itemize}
    \item What do we want to prove?
    \item Lambda calculus (review?)
    \begin{itemize}
        \item Types
        \item Signatures
        \item Syntax
        \item Typing judgments
        \item Denotational semantics
    \end{itemize}
    \item Logical relations
    \begin{itemize}
        \item Types
        \item Signatures
        \item Terms
        \item Fundamental theorem
        \item Application
    \end{itemize}
\end{itemize}
\end{frame}

\begin{frame}
\frametitle{What do we want to prove?}

\pause
\begin{itemize}
    \item Proving something about \emph{all} programs in a language?
    \pause
    \item If we have $+$ and $\times$, how can we prove even in $\Rightarrow$ even out?
    \pause
    \item Logical relations!
    \pause
    \item Also can prove more complicated and interesting theorems.
    \pause
    \item For example:
    \begin{itemize}
        \pause
        \item Termination: do your programs stop?
        \pause
        \item Type safety: do your programs keep going?
        \pause
        \item Optimizations: why can I rewrite my program?
        \pause
        \item Representation independence: internals don't matter if you hide them.
        \pause
        \item Security: show the output doesn't depend on secure information.
    \end{itemize}
\end{itemize}
\end{frame}

\begin{frame}
\frametitle{Types*}

\pause

\[
\alpha, \beta ::= \tau \mid 1 \mid \alpha_1 \times \alpha_2 \mid \alpha \rightarrow \beta
\]

\pause
Where:
\begin{itemize}
    \item $\tau$ is a ground type from a fixed set of symbols, e.g.\ $\{\text{Int}, \text{Bool}, \ldots\},$
    \pause
    \item $1$ is the unit type,
    \pause
    \item $\alpha_1 \times \alpha_2$ is a product type,
    \pause
    \item $\alpha \rightarrow \beta$ is a function type.
\end{itemize}
\end{frame}

\begin{frame}
\frametitle{Signatures}

\pause

A signature $\Sigma = (\Sigma_{\text{const}}, \Sigma_{\text{func}})$ is composed of two sets, namely
\begin{itemize}
    \pause
    \item $\Sigma_{\text{const}}$ whose elements are constant symbols $c : \tau$; and
    \pause
    \item $\Sigma_{\text{func}}$ whose elements are function symbols $f : (\tau_1, \ldots, \tau_n) \to \tau$.
\end{itemize}
\pause
Everything is defined with respect to a signature $\Sigma$.

\medskip

\pause
For example, assume that we have $\text{Int}$ as ground type.
Then we could defined $\Sigma = \left(\left\{\underline{n} : n \in \mathbb{Z},\right\}, \left\{\underline{+}, \underline{\times}\right\}\right)$. 

\end{frame}

\begin{frame}
\frametitle{Syntax*}

\pause

\[
M, N ::= x \mid \langle \rangle \mid c \mid f(M_1, \ldots, M_n) \mid \lambda (x: \alpha). M \mid M \, N \mid \langle M_1, M_2 \rangle \mid \pi_1(M) \mid \pi_2(M)
\]

\pause
Where:
\begin{itemize}
    \item $x$ is a variable from a countably infinite set $\{x, y, z, \ldots \}$,
    \pause
    \item $\langle \rangle$ is the term of type $1$,
    \pause
    \item $c$ is a constant symbol in $\Sigma_{\text{const}}$,
    \pause
    \item $f$ is a function symbol in $\Sigma_{\text{func}}$,
    \pause
    \item $\lambda (x: \alpha). M$ is lambda abstraction with $x$ of type $\alpha$,
    \pause
    \item $M\, N$ is application,
    \pause
    \item $\langle M_1, M_2 \rangle$ is a product,
    \pause
    \item $\pi_1(M)$ and $\pi_2(M)$ are projections.
\end{itemize}
\end{frame}

\begin{frame}
\frametitle{Typing Judgments*}

\pause
\[
\infer{\Gamma \vdash x : \alpha}{(x : \alpha) \in \Gamma}
\qquad
\pause
\infer{\Gamma \vdash \langle \rangle : 1}{}
\pause
\qquad
\infer{\Gamma \vdash c : \tau}{c : \tau \in \Sigma_{\text{const}}}
\]
\pause
\[
\infer{\Gamma \vdash f(M_1, \ldots, M_n) : \tau}{\Gamma \vdash M_1 : \tau_1 & \cdots & \Gamma \vdash M_n : \tau_n & f : (\tau_1, \ldots, \tau_n) \rightarrow \tau \in \Sigma_{\text{func}}}
\]
\pause
\[
\infer{\Gamma \vdash \lambda (x : \alpha). M : \alpha \rightarrow \beta}{\Gamma, x : \alpha \vdash M : \beta}
\pause
\qquad
\infer{\Gamma \vdash M \, N : \beta}{\Gamma \vdash M : \alpha \rightarrow \beta & \Gamma \vdash N : \alpha}
\]
\pause
\[
\infer{\Gamma \vdash \langle M_1, M_2 \rangle : \alpha_1 \times \alpha_2}{\Gamma \vdash M_1 : \alpha_1 & \Gamma \vdash M_2 : \alpha_2}
\pause
\qquad
\infer{\Gamma \vdash \pi_1(M) : \alpha_1}{\Gamma \vdash M : \alpha_1 \times \alpha_2}
\pause
\qquad
\infer{\Gamma \vdash \pi_2(M) : \alpha_2}{\Gamma \vdash M : \alpha_1 \times \alpha_2}
\]
\end{frame}

\begin{frame}
\frametitle{Set-Theoretic Denotational Semantics for Types*}

\pause

Let $\rho$ be a function that assigns a set $\rho(\tau)$ to each ground type $\tau$, e.g., $\rho(\text{Int}) = \mathbb{Z}$, $\rho(\text{Bool}) = \{ \text{true}, \text{false} \}$.

\pause
\medskip

Then to type $\alpha$, we are going to assign a set $\llbracket \alpha \rrbracket_{\rho}$ as follows:
\begin{itemize}
    \pause
    \item $\llbracket \tau \rrbracket_{\rho} = \rho(\tau)$
    \pause
    \item $\llbracket 1 \rrbracket_{\rho} = \{\star\}$
    \pause
    \item $\llbracket \alpha_1 \times \alpha_2 \rrbracket_{\rho} = \llbracket \alpha_1 \rrbracket_{\rho} \times \llbracket \alpha_2 \rrbracket_{\rho}$
    \pause
    \item $\llbracket \alpha \rightarrow \beta \rrbracket_{\rho} = \llbracket \alpha \rrbracket_{\rho} \rightarrow \llbracket \beta \rrbracket_{\rho}$
\end{itemize}
\end{frame}

\begin{frame}
\frametitle{Set-Theoretic Denotational Semantics for Signatures}

\pause

For a fixed $\rho$ assigning ground types to sets, we can give an interpretation $\sigma$  to the constants and functions of a signature $\Sigma = (\Sigma_{\text{const}}, \Sigma_{\text{func}})$:
\begin{itemize}
    \pause
    \item for each $c : \tau \in \Sigma_{\text{const}}$, an element $\sigma(c) \in \rho(\tau)$; and
    \pause
    \item for each $f : (\tau_1, \ldots, \tau_n) \to \tau \in \Sigma_{\text{func}}$, a function $\sigma(f) \in \rho(\tau_1) \times \!\cdots\! \times \rho(\tau_n) \to \rho(\tau)$.
\end{itemize}
\pause
Note that $\sigma(c) \in \llbracket \tau \rrbracket_{\rho}$ and $\sigma(f) \in \llbracket \tau_1 \times \cdots \times \tau_n \rightarrow \tau \rrbracket_{\rho}$.
\end{frame}

\begin{frame}
\frametitle{Set-Theoretic Denotational Semantics for Terms in Context*}

\pause

Fix a fixed $\rho$ and $\sigma$, we can define the meaning of a lambda term.
In a context $\Gamma = x_1 : \alpha_1, \ldots, x_n : \alpha_n$, the denotational semantics of a term $M$ is a function:

\[
\llbracket \Gamma \vdash M : \alpha \rrbracket_{\rho, \sigma} : \llbracket \alpha_1 \times \cdots \times \alpha_n \rrbracket_\rho \to \llbracket \alpha \rrbracket_\rho
\]

For a $\Gamma$ as above, we will write $\gamma$ for an element of $\llbracket \alpha_1 \times \cdots \times \alpha_n \rrbracket_\rho$ and write $\gamma(x_i)$ for the $i^{\text{th}}$ component of the tuple.

\pause
\medskip

We write $\gamma[x \mapsto v]$ to denote the extension of $\gamma$ mapping $x$ to $v$.
E.g.\ for $\Gamma = x : \text{Int}, y : \text{Int}$ if $\{x \mapsto 1, y \mapsto 2\} \in \llbracket \text{Int} \times \text{Int} \rrbracket_\rho$ then
\[
\{x \mapsto 1, y \mapsto 2\}[z \mapsto 3] := \{x \mapsto 1, y \mapsto 2, z \mapsto 3\} \in \llbracket \text{Int} \times \text{Int} \times \text{Int}\rrbracket_\rho
\]
for $\Gamma = x : \text{Int}, y : \text{Int}, z : \text{Int}$.
\end{frame}

\begin{frame}
\frametitle{Recursive Definition of Denotation for Terms*}

\pause

The denotational semantics $\llbracket \Gamma \vdash M : \alpha \rrbracket_{\rho,\sigma} (\gamma)$ is defined recursively as follows:

\begin{itemize}
    \pause
    \item $\llbracket \Gamma \vdash x : \alpha \rrbracket_{\rho,\sigma} (\gamma) = \gamma(x)$
    \pause
    \item $\llbracket \Gamma \vdash \langle \rangle : 1 \rrbracket_{\rho,\sigma} (\gamma) = \star$
    \pause
    \item $\llbracket \Gamma \vdash c : \tau \rrbracket_{\rho,\sigma} (\gamma) = \sigma(c)$
    \pause
    \item $\llbracket \Gamma \vdash f(M_1, \ldots, M_n) : \sigma \rrbracket_{\rho,\sigma} (\gamma) =$ $\phantom{.}\quad\sigma(f)(\llbracket \Gamma \vdash M_1 : \tau_1 \rrbracket_{\rho,\sigma} (\gamma), \ldots, \llbracket \Gamma \vdash M_n : \tau_n \rrbracket_{\rho,\sigma} (\gamma))$
    \pause
    \item $\llbracket \Gamma \vdash \lambda (x : \alpha). M : \alpha \rightarrow \beta \rrbracket_{\rho,\sigma} (\gamma) = \lambda v. \llbracket \Gamma, x:\alpha \vdash M : \beta \rrbracket_{\rho,\sigma} (\gamma[x \mapsto v])$
    \pause
    \item $\llbracket \Gamma \vdash M \, N : \beta \rrbracket_{\rho,\sigma} (\gamma) = \llbracket M \rrbracket_{\rho,\sigma}(\gamma)(\llbracket N \rrbracket_{\rho,\sigma}(\gamma))$
    \pause
    \item $\llbracket \Gamma \vdash \langle M_1, M_2 \rangle : \alpha_1 \times \alpha_2 \rrbracket_{\rho,\sigma} (\gamma) = (\llbracket M_1 \rrbracket_{\rho,\sigma}(\gamma), \llbracket M_2 \rrbracket_{\rho,\sigma}(\gamma))$
    \pause
    \item $\llbracket \Gamma \vdash \pi_1(M) : \alpha_1 \rrbracket_{\rho,\sigma} (\gamma) = \pi_1(\llbracket M \rrbracket_{\rho,\sigma}(\gamma))$
    \pause
    \item $\llbracket \Gamma \vdash \pi_2(M) : \alpha_2 \rrbracket_{\rho,\sigma} (\gamma) = \pi_2(\llbracket M \rrbracket_{\rho,\sigma}(\gamma))$
\end{itemize}
\end{frame}

\begin{frame}
\frametitle{Logical Relations Semantics for Types}

We can view logical relations as a new type of denotational semantics.
We begin with the types.
For each type $\alpha$, we're going to define \emph{two sets}, $P$ and $A$, such that $P \subseteq A$.
The set $P$ is our \emph{predicate}.

\pause
\medskip

Let $\rho$ be a function that assigns a pair of sets $\left(\rho_{\mathcal{P}}(\tau), \rho_{\mathcal{A}}(\tau)\right)$ to each ground type $\tau$ such that $\rho_{\mathcal{P}}(\tau) \subseteq \rho_{\mathcal{A}}(\tau)$, e.g., $\rho(\text{Int}) = \left(\left\{2m : m \in \mathbb{Z}\right\}, \mathbb{Z}\right)$.

\pause
\medskip

Then to type $\alpha$, we are going to assign a pair of sets $\lBrace \alpha \rBrace_{\rho} = \left(\mathcal{P}\lBrace \alpha \rBrace_{\rho}, \mathcal{A}\lBrace \alpha \rBrace_{\rho}\right)$ as follows:

\begin{itemize}
    \pause
    \item $\lBrace \tau \rBrace_{\rho} = \left(\rho_{\mathcal{P}}(\tau), \rho_{\mathcal{A}}(\tau)\right)$
    \pause
    \item $\lBrace 1 \rBrace_{\rho} = \left(\{\star\}, \{\star\}\right)$
    \pause
    \item $\lBrace \alpha_1 \times \alpha_2 \rBrace_{\rho} = \left(\mathcal{P}\lBrace \alpha_1 \rBrace_{\rho} \times \mathcal{P}\lBrace \alpha_2 \rBrace_{\rho}, \mathcal{A}\lBrace \alpha_1 \rBrace_{\rho} \times \mathcal{A}\lBrace \alpha_2 \rBrace_{\rho}\right)$
    \pause
    \item $\lBrace \alpha \rightarrow \beta \rBrace_{\rho} = \left(\left\{ f : \forall x \in \mathcal{P}\lBrace \alpha \rBrace_{\rho}. f(x) \in \mathcal{P}\lBrace \beta \rBrace_{\rho} \right\},\, \mathcal{A}\lBrace \alpha \rBrace_{\rho} \rightarrow \mathcal{A}\lBrace \beta \rBrace_{\rho}\right)$
\end{itemize}

\pause

Note that $\mathcal{A}\lBrace \alpha \rBrace_{\rho} = \llbracket\alpha \rrbracket_{\rho_{\mathcal{A}}}$.

\end{frame}

\begin{frame}
\frametitle{Logical Relations Semantics for Signatures}

\pause

For a fixed $\rho$ assigning ground types to sets, we can give an interpretation $\sigma$  to the constants and functions of a signature $\Sigma = (\Sigma_{\text{const}}, \Sigma_{\text{func}})$:
\begin{itemize}
    \pause
    \item for each $c : \tau \in \Sigma_{\text{const}}$, an element $\sigma(c) \in \rho_{\mathcal{P}}(\tau)$; and
    \pause
    \item for each $f : (\tau_1, \ldots, \tau_n) \to \tau \in \Sigma_{\text{func}}$, a function
    \[\sigma(f) \in \rho_{\mathcal{A}}(\tau_1) \times \cdots \times \rho_{\mathcal{A}}(\tau_n) \to \rho_{\mathcal{A}}(\tau)\]
    such that
    \[(x_1, \ldots, x_n) \in \rho_{\mathcal{P}}(\tau_1) \times \cdots \times \rho_{\mathcal{P}}(\tau_n) \Rightarrow \sigma(f)(x_1, \ldots, x_n) \in \rho_{\mathcal{P}}(\tau).\]
\end{itemize}

\medskip
\pause

Note that $\sigma(c) \in \mathcal{P}\lBrace \tau \rBrace_{\rho}$ and $\sigma(f) \in \mathcal{P}\lBrace \tau_1 \times \cdots \times \tau_n \rightarrow \tau \rBrace_{\rho}$, as well as that $\sigma(c) \in \llbracket \tau \rrbracket_{\rho_{\mathcal{A}}}$ and $\sigma(f) \in \llbracket \tau_1 \times \cdots \times \tau_n \rightarrow \tau \rrbracket_{\rho_{\mathcal{A}}}$.

\medskip
\pause

If we want to forget that $\sigma$ preserves our predicates, we will write $\sigma_{\mathcal{A}}$.

\end{frame}

\begin{frame}
\frametitle{Logical Relations Semantics for Terms in Context}

\pause

Fix a fixed $\rho$ and $\sigma$, we want to define the meaning of a lambda term.
In a context $\Gamma = x_1 : \alpha_1, \ldots, x_n : \alpha_n$, we want an interpretation of type

\[
\lBrace \Gamma \vdash M : \alpha \rBrace_{\rho, \sigma} : \mathcal{A}\lBrace \alpha_1 \times \cdots \times \alpha_n \rBrace_\rho \to \mathcal{A}\lBrace \alpha \rBrace_\rho
\]

such that for all $\gamma \in \mathcal{P}\lBrace \alpha_1 \times \cdots \times \alpha_n \rBrace_\rho$ we have $\lBrace \Gamma \vdash M : \alpha \rBrace_{\rho, \sigma}(\gamma) \in \mathcal{P}\lBrace \alpha \rBrace_\rho$.
I.e., we map values satisfying our predicate to values satisfying our predicate.

\medskip

How do we define this semantics?

\end{frame}

\begin{frame}
\frametitle{Fundamental Theorem of Logical Relations}

Recall that $\mathcal{A}\lBrace \alpha \rBrace_{\rho} = \llbracket\alpha \rrbracket_{\rho_{\mathcal{A}}}$.
Thus, we can define

\[
\lBrace \Gamma \vdash M : \alpha \rBrace_{\rho, \sigma} : \mathcal{A}\lBrace \alpha_1 \times \cdots \times \alpha_n \rBrace_\rho \to \mathcal{A}\lBrace \alpha \rBrace_\rho
\]

as $\lBrace \Gamma \vdash M : \alpha \rBrace_{\rho, \sigma} := \llbracket \Gamma \vdash M : \alpha \rrbracket_{\rho_{\mathcal{A}}, \sigma_{\mathcal{A}}}$ if it actually preserves our predicates.

\medskip
\pause

\begin{theorem}
    Fix $\rho$ and $\sigma$ for logical relations.
    For all $\gamma \in \mathcal{P}\lBrace \alpha_1 \times \cdots \times \alpha_n \rBrace_\rho$ we have $\llbracket \Gamma \vdash M : \alpha \rrbracket_{\rho_{\mathcal{A}}, \sigma_{\mathcal{A}}}(\gamma) \in \mathcal{P}\lBrace \alpha \rBrace_\rho$.
    \begin{proof}
        Induction on the structure of $M$.
    \end{proof}
\end{theorem}

\medskip
\pause

This is known as the \emph{Fundamental Theorem of Logical Relations}, or the \emph{Basic Lemma of Logical Relations}.
Note that we had to choose the interpretation $\sigma$ of our constants and built-in functions to respect $\rho_{\mathcal{P}}$.
\end{frame}

\begin{frame}
\frametitle{Example Application: Preserving Evenness}

\pause

\begin{itemize}
    \item Fix the ground types to be $\{\text{Int}\}$.
    \pause
    \item Fix a signature $\Sigma = \left(\left\{\underline{n} : n \in \mathbb{Z},\, n \text{ even}\right\}, \left\{\underline{+}, \underline{\times}\right\}\right)$, everything using $\text{Int}$.
    \pause
    \item Interpret our ground types with $\rho(\text{Int}) = \left(\left\{2m : m \in \mathbb{Z}\right\}, \mathbb{Z}\right)$.
    \pause
    \item Interpret our signature with $\sigma\left(\underline{n}\right) = n, \sigma\left(\underline{+}\right) = +, \sigma\left(\underline{\times}\right) = \times$ and check that our functions preserve even numbers.
    \pause
    \item Apply the theorem! For example, for all terms $x : \text{Int} \vdash M : \text{Int}$, we have
    \[
    n \in \mathcal{P}\lBrace \text{Int} \rBrace_\rho \Rightarrow \llbracket x : \text{Int} \vdash M : \text{Int} \rrbracket_{\rho_{\mathcal{A}}, \sigma}(n) \in \mathcal{P}\lBrace \text{Int} \rBrace_\rho
    \]
    which is equivalent to
    \[
    n \text{ even} \Rightarrow \llbracket x : \text{Int} \vdash M : \text{Int} \rrbracket_{\rho_{\mathcal{A}}, \sigma}(n) \text{ even}.
    \]
    \vspace{-12pt}
    \pause
    \item Can also show that all polynomials with even coefficients preserve evenness.
    \item Importantly, the theorem also applies to contexts with function types.
    If we have $f : \text{Int} \to \text{Int}$, we are forced to feed in a function from $\mathcal{P}\lBrace \text{Int} \to \text{Int} \rBrace_\rho$, which are exactly even preserving functions!
\end{itemize}

\end{frame}

\begin{frame}
\frametitle{Conclusion}


\begin{itemize}
    \pause
    \item We saw the
    \begin{itemize}
        \item syntax,
        \item typing rules, and
        \item set-theoretic denotational semantics
    \end{itemize}
     of simply typed lambda calculus with products, ground types, constants, and built-in functions.
    \pause
    \item We then extended the denotational semantics to include a predicate at each type, and observed that all of our constructions respected these predicates.
    \pause
    \item The preservation gave use the Fundamental Theorem of Logical Relations, which we then applied to an example.
\end{itemize}

\begin{textblock}{5}[1,0.5](17,4.5)
    \centering
    \vskip4pt
    \pause
    \includegraphics[scale=0.7]{qr-code.eps}
\end{textblock}

\end{frame}

\appendix

\begin{frame}
\frametitle{References}

\printbibliography
\end{frame}

\begin{frame}
\frametitle{Free Variables*}
The function $\text{FV}(M)$ is defined recursively as follows:
\begin{itemize}
    \item $\text{FV}(x) = \{x\}$
    \item $\text{FV}(\langle \rangle) = \emptyset$
    \item $\text{FV}(c) = \emptyset$
    \item $\text{FV}(f(M_1, \ldots, M_n)) = \bigcup_{i=1}^{n} \text{FV}(M_i)$
    \pause
    \item $\text{FV}(\lambda (x: \alpha). M) = \text{FV}(M) \setminus \{x\}$
    \pause
    \item $\text{FV}(M \, N) = \text{FV}(M) \cup \text{FV}(N)$
    \item $\text{FV}(\langle M_1, M_2 \rangle) = \text{FV}(M_1) \cup \text{FV}(M_2)$
    \item $\text{FV}(\pi_1(M)) = \text{FV}(M)$
    \item $\text{FV}(\pi_2(M)) = \text{FV}(M)$
\end{itemize}

\pause
For example, $\text{FV}(\lambda (x: \alpha). y \, x) = \{y\}$.

\end{frame}

\begin{frame}
\frametitle{Capture-avoiding Substitution*}

Substitution of $N$ for $x$ in $M$ in a capture-avoiding way, denoted $M[x := N]$, is defined recursively as follows:

\begin{itemize}
    \item $x[x := N] = N$
    \item $y[x := N] = y$, for $y \neq x$
    \pause
    \item $\langle \rangle[x := N] = \langle \rangle$
    \item $c[x := N] = c$
    \pause
    \item $f(M_1, \ldots, M_n)[x := N] = f(M_1[x := N], \ldots, M_n[x := N]))$
    \item $(M \, P)[x := N] = (M[x := N]) \, (P[x := N])$
    \item $\langle M_1 , M_2 \rangle[x := N] = \langle M_1[x := N], M_2[x := N] \rangle$
    \item $\pi_1(M)[x := N] = \pi_1(M[x := N])$
    \item $\pi_2(M)[x := N] = \pi_2(M[x := N])$
    \item $\ldots$
\end{itemize}
\end{frame}

\begin{frame}
\frametitle{Capture-avoiding Substitution (continued)*}

Most importantly, we have the rule for abstraction:

\begin{itemize}
    \item $(\lambda (y : \alpha). M)[x := N] =
    \begin{cases}
    \lambda (y : \alpha). M[x := N] & \text{if } y \neq x \text{ and } y \notin \text{FV}(N) \\
    \lambda (z : \alpha). M[y := z][x := N] & \text{if } y = x \text{ or } y \in \text{FV}(N)
    \end{cases}$
\end{itemize}
\pause
Here are two examples:
\begin{itemize}
    \item $(\lambda (y : \alpha). x)[x := z] = \lambda (y : \alpha). z$
    \item $(\lambda (y : \alpha). y \, x)[x := y] = \lambda (z : \alpha). z \, y$
\end{itemize}

\end{frame}

\begin{frame}
\frametitle{Equations-in-Context*}

An equation-in-context is expressed as:
\[
\Gamma \vdash M = N : \alpha
\]

The judgments means that in the type context $\Gamma$, the terms $M$ and $N$ are considered equal and both have type $\alpha$.

\medskip

Equations-in-contexts allow us to perform equational reasoning while respecting the types assigned to the variables involved.

\end{frame}

\begin{frame}
\frametitle{Equational Reasoning Rules*}

\[
\infer[\text{(Refl)}]{\Gamma \vdash M = M : \alpha}{\Gamma \vdash M : \alpha}
\]

\[
\infer[\text{(Sym)}]{\Gamma \vdash N = M : \alpha}{\Gamma \vdash M = N : \alpha}
\]

\[
\infer[\text{(Trans)}]{\Gamma \vdash M = P : \alpha}{\Gamma \vdash M = N : \alpha & \Gamma \vdash N = P : \alpha}
\]
\end{frame}

\begin{frame}
\frametitle{Weakening and Substitution Rules*}

\[
\infer[\text{(Weak)}]{\Gamma, x : \beta \vdash M = N : \alpha}{\Gamma \vdash M = N : \alpha}
\]

\[
\infer[\text{(Subs)}]{\Gamma \vdash M[x := P] = N[x := P] : \alpha}{\Gamma, x : \beta \vdash M = N : \alpha & \Gamma \vdash P : \beta}
\]
\end{frame}

\begin{frame}
\frametitle{Rules for Unit and Binary Products*}

\[
\infer[\text{(Unit-Eq)}]{\Gamma \vdash M = \langle\rangle : 1}{\Gamma \vdash M : 1}
\]

\[
\infer[\text{(Proj1)}]{\Gamma \vdash \pi_1(\langle M_1, M_2 \rangle) = M_1 : \alpha_1}{\Gamma \vdash M_1 : \alpha_1 & \Gamma \vdash M_2 : \alpha_2}
\]

\[
\infer[\text{(Proj2)}]{\Gamma \vdash \pi_2(\langle M_1, M_2 \rangle) = M_2 : \alpha_2}{\Gamma \vdash M_1 : \alpha_1 & \Gamma \vdash M_2 : \alpha_2}
\]

\[
\infer[\text{($\eta$-Prod)}]{\Gamma \vdash \langle \pi_1(P), \pi_2(P) \rangle = P : \alpha_1 \times \alpha_2}{\Gamma \vdash P : \alpha_1 \times \alpha_2}
\]
\end{frame}

\begin{frame}
\frametitle{Rules for Functions*}

\[
\infer[\text{($\beta$-Eq)}]{\Gamma \vdash (\lambda (x : \alpha). M) \, N = M[x := N] : \beta}{\Gamma, x : \alpha \vdash M : \beta & \Gamma \vdash N : \alpha}
\]

\[
\infer[\text{($\eta$-Eq)}]{\Gamma \vdash \lambda (x : \alpha). (M \, x) = M : \alpha \rightarrow \beta}{x \not\in \text{FV}(M)}
\]

\[
\infer[\text{($\lambda$-Cong)}]{\Gamma \vdash \lambda (x : \alpha). M = \lambda (x : \alpha). N : \alpha \rightarrow \beta}{\Gamma, x : \alpha \vdash M = N : \beta}
\]

\end{frame}

\begin{frame}
\frametitle{Axioms}
We also want axioms in order to reason about elements of our signature $\Sigma$.
An axiom is a pair of terms $\left(\Gamma \vdash M : \alpha, \Gamma \vdash N : \alpha\right)$ of terms of the same type in the same context.
For a set of axioms $\Omega$, we have the rule
\[
\infer[\text{(Axiom)}]{\Gamma \vdash M = N : \alpha}{\left(\Gamma \vdash M : \alpha, \Gamma \vdash N : \alpha\right) \in \Omega}
\]
Note that it is possible to prove everything equals everything else if you choose your axioms wrong!
\end{frame}

\begin{frame}
\frametitle{Soundness}

\begin{theorem}
    Let $\Omega$ be a set of axioms in a signature $\Sigma$.
    Let $\rho$ and $\sigma$ be assignments such that, for all $\left(\Gamma \vdash M : \alpha, \Gamma \vdash N : \alpha\right) \in \Omega$, we have $\llbracket \Gamma \vdash M : \alpha \rrbracket_{\rho, \sigma} = \llbracket \Gamma \vdash N : \alpha \rrbracket_{\rho, \sigma}$.
    Then, for all valid equations $\Gamma \vdash M = N : \alpha$ we have
    
    \[
        \llbracket \Gamma \vdash M : \alpha \rrbracket_{\rho, \sigma} = \llbracket \Gamma \vdash N : \alpha \rrbracket_{\rho, \sigma}
    \]
    \begin{proof}
        By induction on the proof of $\Gamma \vdash M = N : \alpha$.
    \end{proof}
\end{theorem}

\medskip

Thus, if we respect the axioms, then equivalent terms have equal denotational semantics. This is the minimum we expect from denotational semantics.

\end{frame}

\end{document}
